\specialsection{Заключение}

В дипломной работе была рассмотрена задача рекомендации следующей культуры для сельскохозяйственного поля на основе
истории его использования и базовых признаков поля и региона. Работа была построена как прототип системы поддержки
принятия решений, в которой результат представлен не только в виде
одной наиболее вероятной культуры, но и в виде ранжированного списка кандидатов, сопровождаемого оценкой уверенности и
интерпретационным контекстом.

В ходе работы были подготовлены данные USDA NASS CSB/CDL, выполнено укрупнение культурных классов, построена оконная
выборка по трём предыдущим годам истории поля и зафиксировано воспроизводимое разбиение данных. На этой основе были
реализованы базовые подходы, марковская модель первого порядка и основная модель CatBoost, выбранная в качестве
предиктивного ядра системы.

Экспериментальная оценка показала, что CatBoost устойчиво превосходит простые базовые подходы и Markov-1. При этом
основной результат работы связан с переходом от одновариантного прогноза к рекомендательной постановке: система формирует
короткий список вероятных культур, в котором фактическая культура часто сохраняется среди ближайших альтернатив.

Дополнительно были реализованы слой оценки уверенности, граф знаний и пространственная демонстрация. Слой оценки
уверенности помогает различать более и менее надёжные рекомендации. Граф знаний добавляет к кандидатам предметный
контекст: переходные закономерности, региональную типичность, принадлежность к группам культур, а также поддерживающие и
предупреждающие сигналы. Пространственная демонстрация показывает, как результаты системы могут быть представлены на
уровне отдельных полей в выбранной области карты.

Основной итог работы состоит в том, что задача выбора следующей культуры может быть рассмотрена не только как задача
многоклассового предсказания, но и как задача построения объяснимой рекомендательной системы поддержки принятия решений.
Разработанный прототип позволяет сформировать набор вероятных кандидатов, оценить надёжность рекомендации и получить
интерпретируемый контекст для анализа.